\documentclass[12pt]{article}
\usepackage[utf8]{inputenc}
\usepackage[T1]{fontenc}
\usepackage[margin=1in]{geometry}
\usepackage{setspace}
\usepackage{xcolor}
\usepackage[colorlinks=true, urlcolor=blue!60!black, linkcolor=blue!60!black]{hyperref}
\setlength{\parindent}{2em}
\setlength{\parskip}{1.25em}
\renewcommand{\baselinestretch}{1.0}
\begin{document}
\begin{center}
{\LARGE \textit{Smith v. Doe}, 538 U.S. 84 (2003)\par}
\vspace{0.8em}
{\large Case Brief\par}
\end{center}
\vspace{1em}

\section{Memory Jogger}
Alaska applied its sex offender registration and Internet notification law to offenders convicted before it was enacted; the Court held 6--3 that the law was a civil, nonpunitive regulation under the two-step intent--effects test, so its retroactive application did not violate the Ex Post Facto Clause (\href{https://supreme.justia.com/cases/federal/us/538/84/}{Justia}).

\section{Detailed Case Facts}
In May 1994, Alaska enacted the Alaska Sex Offender Registration Act (ASORA), a ``Megan's Law'' modeled on the nationwide wave of registration statutes incentivized by Congress's Jacob Wetterling Crimes Against Children and Sexually Violent Offender Registration Act (\href{https://www.law.cornell.edu/supct/html/01-729.ZO.html}{Cornell LII Majority}). ASORA imposed two retroactive components on any sex offender or child kidnapper physically present in Alaska:
\begin{itemize}
  \item \textbf{Registration:} Covered individuals had to register with the Department of Corrections before release or with local police within one working day of conviction or entry into the State, submitting their name, aliases, address, employer, vehicle information, photograph, fingerprints, and postconviction treatment history. Single nonaggravated offenders had to verify information annually for 15 years; aggravated or repeat offenders had to verify quarterly for life. Knowing noncompliance triggered criminal prosecution.
  \item \textbf{Notification:} The Alaska Department of Public Safety maintained a central registry. Although fingerprints, driver's license numbers, and medical treatment records stayed confidential, the State published most nonconfidential data---including name, photograph, physical description, home address, employer address, vehicle details, conviction record, and compliance status---on a public Internet website.
\end{itemize}

Respondents John Doe I and John Doe II pleaded \textit{nolo contendere} to aggravated sexual abuse of a minor years before ASORA's enactment, were released from prison in 1990, and completed sex-offender rehabilitation programs (\href{https://www.law.cornell.edu/supct/html/01-729.ZO.html}{Cornell LII Majority}).

\section{Procedural History}
Along with John Doe I's wife, Jane Doe, they sued Alaska officials under 42 U.S.C. \S\ 1983, arguing that retroactive application of ASORA violated the Ex Post Facto Clause. The District Court granted summary judgment for the State. The Ninth Circuit reversed, agreeing that the Alaska Legislature intended a civil regulatory scheme but holding that ASORA's punitive effects overrode that intent (\href{https://www.law.cornell.edu/supct/html/01-729.ZS.html}{Cornell LII Syllabus}). The Supreme Court granted certiorari.

\section{Judicial Votes}
The Supreme Court \textbf{reversed and remanded} the Ninth Circuit by a \textbf{6--3} vote. The case was argued on November 13, 2002, and decided on March 5, 2003 (\href{https://supreme.justia.com/cases/federal/us/538/84/}{Justia}):

\begin{itemize}
    \item \textbf{Majority Opinion:} Justice Kennedy delivered the opinion of the Court, joined by Chief Justice Rehnquist and Justices O'Connor, Scalia, and Thomas.
    \item \textbf{Concurring Opinions:} Justice Thomas filed a concurring opinion. Justice Souter filed an opinion concurring in the judgment.
    \item \textbf{Dissenting Opinions:} Justice Stevens filed a dissenting opinion. Justice Ginsburg filed a dissenting opinion, joined by Justice Breyer.
\end{itemize}

\section{Holding}
\textbf{Questions Presented:}

Whether the retroactive application of a state sex offender registration and public Internet notification statute constitutes retroactive criminal punishment barred by the Ex Post Facto Clause, U.S. Const. art. I, \S\ 10, cl. 1.

\textbf{Holding:} No. In a 6--3 decision, the Supreme Court held that because the Alaska Legislature intended to create a civil, nonpunitive regulatory regime and respondents failed to establish by the ``clearest proof'' that the statute's effects negated that intent, retroactive application of ASORA does not violate the Ex Post Facto Clause (\href{https://www.law.cornell.edu/supct/html/01-729.ZO.html}{Cornell LII Majority}).

\section{Analysis of Opinions}
\subsection{Majority Opinion (Justice Kennedy)}

Justice Kennedy applied the established two-step ``intent--effects'' framework drawn from \textit{Kansas v. Hendricks}, 521 U.S. 346 (1997), and \textit{United States v. Ward}, 448 U.S. 242 (1980), built on the premise of \href{https://supreme.justia.com/cases/federal/us/3/386/}{\textit{Calder v. Bull}}, 3 U.S. (3 Dall.) 386 (1798), that the Ex Post Facto Clause bars only retroactive criminal punishments, not civil regulations.

\subsubsection{Step One: Legislative Intent}

Determining whether a statute is civil or penal is first a question of statutory construction focusing on text and structure. ASORA's text expressly declared that sex offenders pose a high risk of reoffending and identified ``protecting the public from sex offenders'' as the law's ``primary governmental interest'' (\href{https://www.law.cornell.edu/supct/html/01-729.ZO.html}{Cornell LII Majority}). Protecting public safety is a legitimate nonpunitive objective. Although ASORA's registration mandate appeared in Alaska's Code of Criminal Procedure and required judges to notify defendants of registration duties at plea colloquies and sentencing, the Court found those features non-dispositive: the notification provisions resided in the Health, Safety, and Housing Code, the criminal procedure code housed many nonpunitive provisions, and implementation was vested in an administrative agency (the Department of Public Safety) without criminal-process safeguards. Because the legislature intended a civil scheme, only ``the clearest proof'' could override that label.

\subsubsection{Step Two: Punitive Effects (\textit{Mendoza-Martinez} Guideposts)}

The Court evaluated ASORA's effects using the seven non-exhaustive factors from \textit{Kennedy v. Mendoza-Martinez}, 372 U.S. 144, 168--169 (1963), focusing on the five most relevant:
\begin{itemize}
  \item \textbf{History and Traditions:} Sex offender registries are of recent origin and differ fundamentally from colonial shaming, branding, or banishment. Rather than staging face-to-face ridicule or expelling offenders, ASORA disseminates accurate, mostly public criminal-record information passively via an online archive; any stigma is a collateral consequence of conviction.
  \item \textbf{Affirmative Disability or Restraint:} ASORA imposes no physical restraint and is less severe than occupational debarment, which the Court had upheld as civil. Registrants remain free to change jobs or residences without seeking state permission, and the statute did not require in-person updates. Unlike probation or supervised release, noncompliance triggers a separate prosecution rather than revocation of the original sentence.
  \item \textbf{Traditional Aims of Punishment:} Incidental deterrence does not convert a civil regulation into punishment. Differentiating between single nonaggravated offenders (15 years, annual) and aggravated or repeat offenders (lifetime, quarterly) is reasonably tied to recidivism risk rather than retribution.
  \item \textbf{Rational Connection to a Nonpunitive Purpose:} The ``most significant'' factor weighed heavily for Alaska; alerting the public to convicted sex offenders rationally advances public safety even without a perfect fit.
  \item \textbf{Excessiveness:} Relying on legislative findings and citing \textit{McKune v. Lile}, 536 U.S. 24, 34 (2002), for the proposition that sex-offender recidivism risk is ``frightening and high,'' the Court held that a State may regulate convicted offenders categorically without individualized dangerousness hearings. Given population mobility, worldwide Internet access is a reasonable, passive means of informing the public.
\end{itemize}
The final two \textit{Mendoza-Martinez} factors---scienter and whether the underlying conduct is already a crime---carried little weight because registration is triggered by past conviction as a baseline for assessing recidivism, not by a new criminal act.

\subsection{Concurring Opinion (Justice Thomas)}

Justice Thomas joined the majority in full but wrote separately to insist that Ex Post Facto analysis must examine only ``the statute on its face,'' not executive implementation (\href{https://www.law.cornell.edu/supct/html/01-729.ZC.html}{Cornell LII Concurrence}). Because ASORA did not mandate Internet publication, the Court strayed by evaluating online dissemination.

\subsection{Opinion Concurring in the Judgment (Justice Souter)}

Justice Souter rejected the ``clearest proof'' presumption because ASORA omitted an express ``civil'' label and split its codification between civil and criminal titles (\href{https://www.law.cornell.edu/supct/html/01-729.ZC1.html}{Cornell LII Concurrence}). Viewing the civil and punitive indicators as in ``rough equipoise,'' Souter concurred solely on the general presumption of constitutionality afforded state laws.

\subsection{Dissenting Opinion (Justice Stevens)}

Justice Stevens argued that ASORA severely impairs protected liberty and is indisputably punitive because a criminal conviction is both a necessary and sufficient condition for its burdens (\href{https://www.law.cornell.edu/supct/html/01-729.ZD.html}{Cornell LII Dissent}). He proposed a bright-line rule: any sanction imposed on everyone convicted of a crime, on no one else, and severely impairing liberty is punishment. (He concurred in the judgment in companion case \textit{Connecticut Dept. of Public Safety v. Doe}, reasoning that registries are permissible \textit{prospective} punishment.)

\subsection{Dissenting Opinion (Justice Ginsburg)}

Justice Ginsburg, joined by Justice Breyer, agreed with Souter that legislative intent was ambiguous and evaluated the \textit{Mendoza-Martinez} factors neutrally (\href{https://www.law.cornell.edu/supct/html/01-729.ZD1.html}{Cornell LII Dissent}). She found ASORA punitive in effect---resembling supervised release and public shaming---and fatally excessive because it imposed lifetime quarterly reporting without individualized risk assessment or any mechanism for rehabilitated offenders like John Doe I to exit the registry.

\subsection{Subsequent Developments}

\begin{itemize}

\item \textbf{2026 Supreme Court Application (\textit{Ellingburg v. United States}):} On January 20, 2026, a unanimous Court in \href{https://www.law.cornell.edu/supremecourt/text/24-482}{\textit{Ellingburg v. United States}}, 607 U.S. \_\_\_ (2026), applied \textit{Smith}'s Step One text-and-structure test to hold that restitution under the Mandatory Victims Restitution Act (MVRA) is criminal punishment subject to the Ex Post Facto Clause. Writing for the Court, Justice Kavanaugh distinguished \textit{Smith} because Alaska adopted ``distinctly civil procedures,'' whereas MVRA restitution is labeled a ``penalty'' in the criminal code, imposed at sentencing, and enforced by the United States. Justice Thomas (joined by Justice Gorsuch) concurred, urging a return to \textit{Calder v. Bull}'s originalist framework.

\end{itemize}

\section{Examples: Future Applications}
\begin{enumerate}

\item \textbf{Same-Side (Civil Regulation; No Ex Post Facto Violation): A Retroactive Financial-Fraud Registry.} Congress enacts the ``Investor Protection Act,'' codified in the civil securities title, requiring anyone convicted of securities fraud since 1995 to report all bank accounts quarterly for life and publishing their home address, photo, and employer on a searchable SEC ``Fraudster Watch'' website, with no off-ramp for rehabilitation. Under \textit{Smith}, the government argues civil intent (express investor-protection purpose, civil codification, administrative enforcement) and nonpunitive effects (no physical restraint, passive online archive of public convictions, rational fit). Challengers argue under \textit{Snyder} and Justice Ginsburg's dissent that lifetime reporting of private financial accounts and home addresses without individualized risk assessment is excessive and stigmatizing, though \textit{Smith}'s ``clearest proof'' standard heavily favors the government unless geographic or occupational bans are added.

\item \textbf{Same-Side (Civil Regulation; No Ex Post Facto Violation): Civil Commitment of a Dangerous Offender.} As an inmate convicted of child sexual abuse nears release, the State petitions under a newly enacted statute to commit him to a secure treatment facility as a sexually violent predator, relying partly on conduct that predates the law. The statute is placed in the probate code, requires proof of a mental abnormality making the person dangerous, and provides for release once he is no longer dangerous. Under the intent--effects test that \textit{Smith} drew from \textit{Kansas v.\ Hendricks}, 521 U.S.\ 346 (1997), the legislature's civil purpose controls unless the clearest proof shows punitive effect, and \textit{Hendricks} upheld materially identical commitment against an Ex Post Facto challenge.

\item \textbf{Opposite-Side (Punishment): Residency Bans Enforced by the Sentencing Court.} Suppose instead Congress amends federal criminal statutes to impose a retroactive ``Public Safety Surcharge'' and a 1,500-foot residential exclusion zone around schools and parks on all pre-enactment offenders, enforced by the sentencing court at a criminal modification hearing with immediate re-incarceration for any technical violation. Unlike \textit{Smith}, this statute fails at \textit{both} steps. Under Step One and \href{https://www.law.cornell.edu/supremecourt/text/24-482}{\textit{Ellingburg v. United States}} (2026), codification in the criminal code, imposition by a sentencing court, and enforcement through criminal resentencing reveal punitive legislative intent, ending the inquiry. Even if Congress labeled it ``civil,'' under Step Two and \textit{Snyder} / \href{https://law.justia.com/cases/federal/appellate-courts/ca6/23-5248/23-5248-2024-05-15.html}{\textit{Doe v. Lee}} (6th Cir. 2024), retroactive residential banishment and revocation-style enforcement impose affirmative physical restraints historically regarded as punishment, supplying the ``clearest proof'' needed to invalidate the law under the Ex Post Facto Clause.

\item \textbf{Opposite-Side (Punishment): Applying Harsher Sentencing Guidelines to an Earlier Crime.} A defendant commits bank fraud in 2005 and is sentenced in 2010 under amended federal Sentencing Guidelines that raise his recommended range well above the range in effect when he acted. No one disputes that a prison sentence is punishment, so there is no civil label for \textit{Smith}'s intent--effects inquiry to test; the only question is whether the change increases the punishment attached to his crime. In \textit{Peugh v.\ United States}, 569 U.S.\ 530 (2013), the Court held that sentencing a defendant under Guidelines promulgated after the offense, when they yield a higher range, violates the Ex Post Facto Clause, even though the Guidelines are advisory.

\item \textbf{Fence-Sitter (Close Call): Lifetime GPS Monitoring Run by the Corrections Department.} A State requires every registrant convicted of an aggravated sex offense, including those convicted before the law took effect, to wear a GPS ankle monitor for life, with the corrections department tracking his location around the clock. The statute declares a public-safety purpose and sits in the civil code, so Step One favors the State. At Step Two, however, the device is a continuous physical burden supervised by corrections officers, resembling the probation and supervised release that \textit{Smith} distinguished because Alaska's registrants could move and work without supervision. Whether that restraint supplies the clearest proof of punishment, or is a reasonable means of protecting the public, has divided courts and is uncertain.

\end{enumerate}

\section{Critique}
\textbf{The ``Clearest Proof'' Deferential Trap:} Critics echo Justices Souter and Ginsburg: once a legislature utters ``public safety'' in a preamble, \textit{Smith}'s ``clearest proof'' threshold makes Step Two nearly insurmountable, inviting legislative relabeling of harsh collateral penalties as ``civil.''

\textbf{Empirical Fragility:} \textit{Smith} rested heavily on the claim that sex-offender recidivism is ``frightening and high''---a phrase traced to a non-peer-reviewed \textit{Psychology Today} article later disavowed by its author, as subsequent courts (including \textit{Snyder}) and criminologists have noted (\href{https://www.narsol.org/scotus0125/}{NARSOL}).

\textbf{Dated Internet Analogy:} Treating a searchable, algorithmically indexed public registry as equivalent to ``a visit to an official archive'' understates the permanent digital stigma, automated background screening, and private ostracism of the modern Internet era.

\section{Key Quotations}
\begin{quote}
``If the intention of the legislature was to impose punishment, that ends the inquiry. If, however, the intention was to enact a regulatory scheme that is civil and nonpunitive, we must further examine whether the statutory scheme is `so punitive either in purpose or effect as to negate [the State's] intention' to deem it `civil.'\,''
\end{quote}
\noindent --- Justice Kennedy, \href{https://www.law.cornell.edu/supct/html/01-729.ZO.html}{538 U.S. at 92}.

\begin{quote}
``The process is more analogous to a visit to an official archive of criminal records than it is to a scheme forcing an offender to appear in public with some visible badge of past criminality. The Internet makes the document search more efficient, cost effective, and convenient for Alaska's citizenry.''
\end{quote}
\noindent --- Justice Kennedy, \href{https://www.law.cornell.edu/supct/html/01-729.ZO.html}{538 U.S. at 99}.

\begin{quote}
``The excessiveness inquiry of our ex post facto jurisprudence is not an exercise in determining whether the legislature has made the best choice possible to address the problem it seeks to remedy. The question is whether the regulatory means chosen are reasonable in light of the nonpunitive objective.''
\end{quote}
\noindent --- Justice Kennedy, \href{https://www.law.cornell.edu/supct/html/01-729.ZO.html}{538 U.S. at 105}.

\begin{quote}
``In my opinion, a sanction that (1) is imposed on everyone who commits a criminal offense, (2) is not imposed on anyone else, and (3) severely impairs a person's liberty is punishment.''
\end{quote}
\noindent --- Justice Stevens, dissenting, \href{https://www.law.cornell.edu/supct/html/01-729.ZD.html}{538 U.S. at 113}.

\begin{quote}
``And meriting heaviest weight in my judgment, the Act makes no provision whatever for the possibility of rehabilitation: Offenders cannot shorten their registration or notification period, even on the clearest demonstration of rehabilitation or conclusive proof of physical incapacitation.''
\end{quote}
\noindent --- Justice Ginsburg, dissenting, \href{https://www.law.cornell.edu/supct/html/01-729.ZD1.html}{538 U.S. at 117}.


\end{document}
